\chapter{Crabs (Pubic Lice)}

\section*{Definition and Overview}
Crabs, properly called pubic lice, are tiny insects that live in coarse body hair, most commonly the pubic hair. They are a common sexually transmissible infection (STI), spread mainly through close physical or sexual contact, and can also be passed on through shared bedding, towels, or clothing. The main symptom is intense itching in the affected area. Pubic lice do not carry other diseases, and they are easily treated with creams or lotions available from a pharmacy. Both sexual partners need treatment, and clothing and bedding should be washed to prevent re-infestation.

\section*{Epidemiology}
Pubic lice are one of the more common STIs, although their incidence has fallen in recent decades, possibly because of changes in pubic hair grooming. They affect men and women of all ages, with the highest rates in people in their teens and twenties. Pubic lice are spread by close contact, and having one STI increases the risk of others, so testing for other infections is recommended. They are not related to poor hygiene, and anyone can get them.

\section*{Aetiology and Pathophysiology}
Pubic lice are caused by the louse \textit{Pthirus pubis}, a small crab-shaped insect.

\begin{itemize}
    \item \textbf{The louse:} \textit{Pthirus pubis} is a wingless insect that attaches to coarse body hair, mainly in the pubic area
    \item \textbf{Transmission:} spread by close physical or sexual contact with an infected person
    \item \textbf{Fomites:} can survive briefly on bedding, towels, and clothing
    \item \textbf{Nits:} the lice lay eggs (nits) attached to the base of hairs
    \item \textbf{Feeding:} lice bite the skin and feed on blood, causing itching
    \item \textbf{Survival:} lice cannot survive away from the body for more than about one to two days
\end{itemize}

\section*{Clinical Features}
Symptoms usually appear about five days to several weeks after exposure.

\begin{itemize}
    \item Intense itching in the pubic area, worse at night
    \item Visible lice or nits (small pale eggs) attached to the base of pubic hairs
    \item Tiny red or blue-grey spots on the skin where the lice have fed
    \item Inflammation or secondary infection from scratching
    \item Lice may also be found in armpit, chest, beard, or eyelash hair
    \item A feeling of something moving in the hair
\end{itemize}

\section*{Differential Diagnosis}
Other conditions cause itching in the genital area and should be considered.

\begin{itemize}
    \item Scabies, caused by a mite that burrows into the skin
    \item Fungal infection (tinea cruris, "jock itch")
    \item Contact dermatitis from soaps or products
    \item Eczema or psoriasis
    \item Bacterial vaginosis, thrush, or other genital infections
    \item Body lice, which live in clothing and bedding
    \item Other sexually transmitted infections that can cause genital irritation
\end{itemize}

\section*{When to Seek Medical Care}
Pubic lice can usually be treated at home with products from the pharmacy, but medical review is sensible if symptoms are severe, the diagnosis is uncertain, treatment fails, or there is a chance of another STI.

\textbf{Call 000 (or local emergency services) for:}
\begin{itemize}
    \item Severe allergic reaction with swelling of the face, lips, or throat, or difficulty breathing after applying a treatment product
    \item Signs of a serious skin infection with spreading redness, pain, and fever
    \item Intense itching with a widespread rash and fever, which may indicate another condition
\end{itemize}

Seek medical care if the itching is severe, the skin becomes infected, symptoms persist after treatment, the lice affect the eyelashes, or you are concerned about other sexually transmitted infections.

\section*{Diagnosis and Investigations}
Diagnosis is usually made by inspection.

\begin{itemize}
    \item \textbf{Examination:} the pubic and other coarse hair is inspected for lice and nits
    \item \textbf{Identification:} lice or eggs may be removed with fine tweezers or a comb and identified
    \item \textbf{STI screening:} testing for chlamydia, gonorrhoea, HIV, and syphilis is recommended, as risk factors overlap
    \item \textbf{Partner notification:} sexual partners from the previous few months should be informed and treated
\end{itemize}

\section*{Management}
Treatment is simple and effective.

\begin{itemize}
    \item \textbf{Permethrin cream or lotion:} applied to the affected hair and washed off after the recommended time
    \item \textbf{Other treatments:} malathion or benzyl benzoate lotions as alternatives
    \item \textbf{Repeated application:} a second application after seven to ten days to kill newly hatched lice
    \item \textbf{Combing:} removal of nits with a fine-toothed comb
    \item \textbf{Treating partners:} all sexual partners should be treated at the same time to prevent re-infection
    \item \textbf{Washing:} bedding, towels, and clothing washed in hot water or dried in a hot dryer
    \item \textbf{Avoiding sex:} until treatment is complete
    \item \textbf{Eyelash treatment:} ointment or petroleum jelly applied with medical advice, never insecticide on the eyes
\end{itemize}

\section*{Complications}
\begin{itemize}
    \item Secondary bacterial infection of scratched skin
    \item Persistent infestation from incomplete treatment or re-infection from partners
    \item Eyelid irritation when lice affect the eyelashes
    \item The psychological and social distress associated with an STI
    \item The risk of other STIs, which is why testing is advised
\end{itemize}

\section*{Prognosis and Outcomes}
Pubic lice are easily cured with correct treatment, and symptoms settle within a few days once the lice are killed. Itching may persist for a week or two as the skin heals. Re-infestation is common if partners are not treated or bedding is not washed. There are no long-term health consequences of pubic lice themselves, and full recovery is expected with treatment.

\section*{Prevention and Self-Care}
\begin{itemize}
    \item Practise safer sex; condoms do not fully prevent pubic lice, which spread by skin and hair contact
    \item Avoid sharing towels, bedding, and clothing with others
    \item Treat all sexual partners at the same time
    \item Wash bedding and clothing in hot water after treatment
    \item Complete the full treatment course, including the second application
    \item Get tested for other STIs if you have had unprotected sex
    \item Revisit a doctor if symptoms persist after treatment
\end{itemize}

\section*{Sources and Further Reading}
The following reputable sources provide further information for healthcare students and patients seeking reliable, current advice about pubic lice.

\begin{itemize}
    \item Healthdirect Australia. \textit{Pubic lice (crabs).} https://www.healthdirect.gov.au/
    \item NSW Health. \textit{Pubic lice fact sheet.} https://www.health.nsw.gov.au/
    \item NHS. \textit{Pubic lice.} https://www.nhs.uk/conditions/pubic-lice/
    \item Centers for Disease Control and Prevention. \textit{Pubic "crab" lice.} https://www.cdc.gov/
    \item Sexual Health Infolink. \textit{Pubic lice.} https://shil.tv/
    \item MedlinePlus. \textit{Pubic lice.} https://medlineplus.gov/ency/article/000841.htm
\end{itemize}
